% ============================================================================
% Tests report: tracking quality.
% ============================================================================
\section{Calidad del seguimiento}\label{sec:calidad}

Los tres ensayos se siguieron completos, con los 22 robots en todos los cuadros: no quedó ninguna
posición predicha ni perdida, y el \SI{99.9}{\percent} de los puntos se detectó con el umbral
estricto (\cref{tab:calidad}). Los errores se corrigieron a mano con anclas; las advertencias se
revisaron visualmente y, por ser pocas y breves, se dejaron como están. Parte de ellas se debe a un
obstáculo que cruza la escena durante un intervalo corto del ensayo, que no afecta al resto de la
medición.

\begin{table}[H]
  \centering
  \caption{Estados de las posiciones. Un \emph{episodio} es un grupo de cuadros seguidos (separados por
    menos de \Cuadros{3}) con al menos un robot no detectado con el umbral estricto.}
  \label{tab:calidad}
  \small
  \begin{tabular}{@{}lrrr@{}}
    \toprule
    & \Cmp{Lab}{A} & \Cmp{Lab}{B} & \Cmp{Lab}{C} \\
    \midrule
    Puntos (robots $\times$ cuadros) & \num{265518} & \num{217382} & \num{242264} \\
    Detectados con umbral estricto & \num{265298} & \num{217149} & \num{241966} \\
    Detectados con umbral relajado & \num{\Cmp{Rel}{A}} & \num{\Cmp{Rel}{B}} & \num{\Cmp{Rel}{C}} \\
    Interpolados & \num{\Cmp{Int}{A}} & \num{\Cmp{Int}{B}} & \num{\Cmp{Int}{C}} \\
    Anclas manuales & \num{\Cmp{Man}{A}} & \num{\Cmp{Man}{B}} & \num{\Cmp{Man}{C}} \\
    Predichos o perdidos & \num{\Cmp{Pre}{A}} & \num{\Cmp{Pre}{B}} & \num{\Cmp{Pre}{C}} \\
    \midrule
    Episodios con advertencias & 46 & 32 & 16 \\
    Cuadros afectados & \SI{1.81}{\percent} & \SI{2.22}{\percent} & \SI{2.69}{\percent} \\
    Episodio más largo & \SI{16}{\second} & \SI{11}{\second} & \SI{52}{\second} \\
    Robots involucrados & 11 & 7 & 5 \\
    Rotación medida (no interpolada) & \SI{100}{\percent} & \SI{99.99}{\percent} & \SI{100}{\percent} \\
    \bottomrule
  \end{tabular}
\end{table}

Las tres anclas manuales están en los últimos cuadros de \Cmp{Lab}{B} (robot 17, cuadros 9874 y
9875) y de \Cmp{Lab}{C} (robot 19, cuadro 11007). El episodio más largo, en \Cmp{Lab}{C} entre
\SI{19.2}{\minute} y \SI{20.1}{\minute}, es un único robot (el 20) detectado con umbral relajado,
con coberturas del anillo de \num{0.73} a \num{0.78}: su anillo se veía parcialmente tapado, pero la
posición es correcta. Como todos los puntos quedan con su estado en la columna \texttt{status},
pueden excluirse en cualquier análisis sensible; con menos del \SI{0.13}{\percent} de puntos
afectados, incluirlos no cambia ningún resultado de este informe.
